\textbf{Phần 1. Trắc nghiệm lựa chọn}
\setcounterpageref{ex}{0}

\begin{ex} %[3D2V23]
Tính giá trị của biểu thức: $A = 24 \times 3 - 35 : 5$.
\choice
{$60$}
{$70$}
{\True $65$}
{$58$}
\loigiai{
Thực hiện nhân và chia trước: $24 \times 3 = 72$; $35 : 5 = 7$.\
Biểu thức có giá trị là: $72 - 7 = 65$.
\textbf{Đáp án của câu hỏi là C}
}
\end{ex}

\begin{ex} %[3D2V23]
Một cửa hàng nhập về 4 thùng bánh, mỗi thùng có 18 hộp bánh. Cửa hàng đã bán được 25 hộp bánh. Hỏi cửa hàng còn lại bao nhiêu hộp bánh?
\choice
{\True $47$ hộp}
{$52$ hộp}
{$45$ hộp}
{$42$ hộp}
\loigiai{
Tổng số hộp bánh nhập về là: $18 \times 4 = 72$ (hộp).\
Số hộp bánh còn lại là: $72 - 25 = 47$ (hộp).
\textbf{Đáp án của câu hỏi là A}
}
\end{ex}

\begin{ex} %[3D2V23]
Từ ba chữ số $2, 4, 6$, người ta lập một số có hai chữ số và một số có một chữ số (mỗi chữ số chỉ dùng một lần) để được tích lớn nhất có thể. Tích lớn nhất đó là
\choice
{$240$}
{$264$}
{$288$}
{\True $252$}
\loigiai{
Để tích lớn nhất, chữ số lớn nhất $6$ phải là thừa số có một chữ số, số có hai chữ số là $42$. Tích là: $42 \times 6 = 252$.
\textbf{Đáp án của câu hỏi là D}
}
\end{ex}

\begin{ex} %[3D2V23]
Một hình tam giác có 3 cạnh bằng nhau, mỗi cạnh dài $28\text{ cm}$. Chu vi của hình tam giác đó là
\choice
{$76\text{ cm}$}
{\True $84\text{ cm}$}
{$80\text{ cm}$}
{$88\text{ cm}$}
\loigiai{
Chu vi của tam giác là: $28 \times 3 = 84\text{ (cm)}$.
\textbf{Đáp án của câu hỏi là B}
}
\end{ex}

\begin{ex} %[3D2V23]
Tìm $y$, biết: $y : 4 = 16 + 8$.
\choice
{$88$}
{$104$}
{\True $96$}
{$92$}
\loigiai{
Ta có: $y : 4 = 24 \Rightarrow y = 24 \times 4 = 96$.
\textbf{Đáp án của câu hỏi là C}
}
\end{ex}

\begin{ex} %[3D2V23]
Mỗi cuộn len đan được 3 chiếc găng tay len. Hỏi có 15 cuộn len như vậy thì đan được bao nhiêu đôi găng tay len (biết 1 đôi gồm 2 chiếc, nếu còn lẻ không tính)?
\choice
{\True $22$ đôi găng tay}
{$45$ đôi găng tay}
{$23$ đôi găng tay}
{$21$ đôi găng tay}
\loigiai{
Tổng số găng tay len đan được là: $15 \times 3 = 45$ (chiếc).\
Số đôi găng tay đan được là: $45 : 2 = 22$ (dư 1 chiếc). Vậy đan được 22 đôi.
\textbf{Đáp án của câu hỏi là A}
}
\end{ex}

\begin{ex} %[3D2V23]
Lớp 3A có 3 tổ, mỗi tổ có 12 bạn. Lớp 3B có nhiều hơn lớp 3A là 5 bạn. Hỏi lớp 3B có bao nhiêu bạn học sinh?
\choice
{$36$ bạn}
{$40$ bạn}
{$42$ bạn}
{\True $41$ bạn}
\loigiai{
Số bạn của lớp 3A là: $12 \times 3 = 36$ (bạn).\
Số bạn của lớp 3B là: $36 + 5 = 41$ (bạn).
\textbf{Đáp án của câu hỏi là D}
}
\end{ex}

\begin{ex} %[3D2V23]
Một mảnh vườn hình chữ nhật có chiều rộng $14\text{ m}$, chiều dài gấp 3 lần chiều rộng. Chu vi của mảnh vườn đó là
\choice
{$84\text{ m}$}
{\True $112\text{ m}$}
{$98\text{ m}$}
{$120\text{ m}$}
\loigiai{
Chiều dài mảnh vườn là: $14 \times 3 = 42\text{ (m)}$.\
Chu vi mảnh vườn là: $(42 + 14) \times 2 = 56 \times 2 = 112\text{ (m)}$.
\textbf{Đáp án của câu hỏi là B}
}
\end{ex}

\begin{ex} %[3D2V23]
Tính nhanh giá trị của biểu thức: $16 \times 3 + 16 \times 2 + 16$.
\choice
{$80$}
{$112$}
{\True $96$}
{$90$}
\loigiai{
Viết lại: $16 \times 3 + 16 \times 2 + 16 \times 1 = 16 \times (3 + 2 + 1) = 16 \times 6 = 96$.
\textbf{Đáp án của câu hỏi là C}
}
\end{ex}

\begin{ex} %[3D2V23]
Tìm một số biết rằng lấy số đó chia cho 3 rồi trừ đi 15 thì được kết quả là 12. Số đó là
\choice
{\True $81$}
{$72$}
{$90$}
{$69$}
\loigiai{
Số đó chia cho 3 là: $12 + 15 = 27$. Số cần tìm là: $27 \times 3 = 81$.
\textbf{Đáp án của câu hỏi là A}
}
\end{ex}

\begin{ex} %[3D2V23]
Bác nông dân nuôi 4 lồng gà, mỗi lồng có 16 con gà. Bác đã bán đi một nửa số gà đó. Hỏi bác nông dân còn lại bao nhiêu con gà?
\choice
{$64$ con gà}
{$28$ con gà}
{$36$ con gà}
{\True $32$ con gà}
\loigiai{
Tổng số gà là: $16 \times 4 = 64$ (con). Bác còn lại một nửa là: $64 : 2 = 32$ (con gà).
\textbf{Đáp án của câu hỏi là D}
}
\end{ex}

\begin{ex} %[3D2V23]
Trong một phép nhân, nếu giữ nguyên thừa số thứ hai và tăng thừa số thứ nhất thêm 5 đơn vị thì tích tăng thêm 35 đơn vị. Thừa số thứ hai là
\choice
{$5$}
{\True $7$}
{$8$}
{$6$}
\loigiai{
Khi thừa số thứ nhất tăng 5 đơn vị thì tích tăng thêm 5 lần thừa số thứ hai.\
Thừa số thứ hai là: $35 : 5 = 7$.
\textbf{Đáp án của câu hỏi là B}
}
\end{ex}

\textbf{Phần 2. Trắc nghiệm đúng sai}
\setcounterpageref{ex}{0}

\begin{ex} %[3D2V23]
Xét tính đúng sai của các phép tính nhân số có hai chữ số với số có một chữ số:
\choiceTF
{\True Khi nhân số có hai chữ số với số có một chữ số, ta thực hiện nhân lần lượt từ phải sang trái.}
{Tích của số có hai chữ số với số có một chữ số luôn là số có hai chữ số.}
{\True Kết quả của $25 \times 4$ bằng $100$.}
{Phép nhân $16 \times 5$ có kết quả là $70$.}
\loigiai{
\begin{itemize}
    \item Ý a) Đúng: Quy tắc đặt tính nhân.
    \item Ý b) Sai: Sai vì ví dụ $50 \times 3 = 150$ là số có 3 chữ số.
    \item Ý c) Đúng: $25 \times 4 = 100$.
    \item Ý d) Sai: Sai vì $16 \times 5 = 80$.
\end{itemize}
\begin{itemize}
    \item a) \textbf{Đúng}.
    \item b) \textbf{Sai}.
    \item c) \textbf{Đúng}.
    \item d) \textbf{Sai}.
\end{itemize}
}
\end{ex}

\begin{ex} %[3D2V23]
Một tổ sản xuất có 4 công nhân, mỗi ngày mỗi công nhân đóng gói được 24 hộp quà:
\choiceTF
{\True Mỗi ngày cả tổ đóng gói được 96 hộp quà.}
{\True Trong 2 ngày, mỗi công nhân đóng gói được 48 hộp quà.}
{Sau 3 ngày cả tổ đóng gói được đúng 200 hộp quà.}
{\True Nếu có thêm 1 công nhân cùng năng suất thì mỗi ngày cả tổ đóng gói được 120 hộp quà.}
\loigiai{
\begin{itemize}
    \item Ý a) Đúng: $24 \times 4 = 96$ hộp quà.
    \item Ý b) Đúng: $24 \times 2 = 48$ hộp quà.
    \item Ý c) Sai: Sai vì $96 \times 3 = 288$ hộp quà.
    \item Ý d) Đúng: $24 \times 5 = 120$ hộp quà.
\end{itemize}
\begin{itemize}
    \item a) \textbf{Đúng}.
    \item b) \textbf{Đúng}.
    \item c) \textbf{Sai}.
    \item d) \textbf{Đúng}.
\end{itemize}
}
\end{ex}

\textbf{Phần 3. Trả lời ngắn}
\setcounterpageref{ex}{0}

\begin{ex} %[3D2V23]
Tính giá trị của biểu thức: $18 \times 5 - 42$.
\shortans{48}
\loigiai{
\textbf{Đáp số:} $48$\\
$18 \times 5 = 90$; $90 - 42 = 48$.
}
\end{ex}

\begin{ex} %[3D2V23]
Một cuộn dây điện dài $15\text{ m}$. Người ta dùng 6 cuộn dây điện như thế. Hỏi có tất cả bao nhiêu mét dây điện?
\shortans{90}
\loigiai{
\textbf{Đáp số:} $90$\\
$15 \times 6 = 90\text{ (m)}$.
}
\end{ex}

\begin{ex} %[3D2V23]
Mẹ mua 4 vỉ sữa chua, mỗi vỉ có 4 hộp. Sau đó mẹ mua thêm 2 vỉ sữa chua loại 6 hộp. Hỏi mẹ mua tất cả bao nhiêu hộp sữa chua?
\shortans{28}
\loigiai{
\textbf{Đáp số:} $28$\\
Số hộp là: $4 \times 4 + 2 \times 6 = 16 + 12 = 28$ (hộp).
}
\end{ex}

\begin{ex} %[3D2V23]
Một hình vuông có cạnh dài $23\text{ cm}$. Tính chu vi của hình vuông đó theo xăng-ti-mét.
\shortans{92}
\loigiai{
\textbf{Đáp số:} $92$\\
Chu vi là: $23 \times 4 = 92\text{ (cm)}$.
}
\end{ex}

\begin{ex} %[3D2V23]
Tìm giá trị của $x$, biết: $x : 3 = 14 + 18$.
\shortans{96}
\loigiai{
\textbf{Đáp số:} $96$\\
$x : 3 = 32 \Rightarrow x = 32 \times 3 = 96$.
}
\end{ex}

\begin{ex} %[3D2V23]
Có 3 xe chở hàng, mỗi xe chở 26 bao xi măng. Người ta đã dỡ xuống 38 bao xi măng. Hỏi trên các xe còn lại bao nhiêu bao xi măng?
\shortans{40}
\loigiai{
\textbf{Đáp số:} $40$\\
Tổng số bao xi măng: $26 \times 3 = 78$ bao. Còn lại: $78 - 38 = 40$ bao.
}
\end{ex}
